\scanpage{169}
\mspage{383}

for the moment to ignore this possibility of normalization -- which brings out more clearly, indeed, the crucial character, in all these $l$-adic calculations, of the role played by the four quantities $c,d,t,u\in\widehat{\mathbb Z}$.

\scanpage{170}
\mspage{384}

\uline{45} \uline{The homomorphism} $\widetilde M_{0,3}\longrightarrow\mathrm{Gl}(2,\widehat{\mathbb Z})'$, \uline{construction of the extension} \deleted{\illegible{1}} $\widetilde N_{1,1}$ \uline{of} $\widetilde\Gamma$ \uline{by} $\mathrm{Sl}(2,\mathbb Z)^\wedge$.

We denote by $\widetilde M_{0,3}$ the group of \deleted{outer} loop automorphisms of $\unsure{\widehat{\mathcal G}_{0,3}^{+}}$ which \deleted{\unsure{commute}} \inserted{induce} up to conjugacy \deleted{with the action of $\unsure{\mathfrak S_3|\mathcal G_3}$, \unsure{which is}} \inserted{on the quotient $\unsure{\mathfrak S_3|\mathcal G_3}\simeq\widehat{\mathcal G}_{0,3}^{+}/\widehat\Pi_{0,3}$ the trivial action -- we have} thus a \inserted{structure of an} extension
\[
(1)\qquad 1\longrightarrow\unsure{\widehat{\mathcal G}_{0,3}^{+}}\longrightarrow\widetilde M_{0,3}\longrightarrow\widetilde N_{0,3}\longrightarrow 1,
\]
\margin{(We shall denote by $\widetilde M_{0,3}^{!}$ the subgroup consisting of the \unsure{automorphisms} which induce \unsure{strict\textsuperscript{ly}} \unsure{the ident\textsuperscript{ity}} in $\unsure{\mathfrak S_3|\mathcal G_3}$ -- it corresponds to $\widetilde N_{0,3}^{!}$ \unsure{-- automorphisms} that are outer loop automorphisms fixing the loop classes (and \unsure{\deleted{induce}|induce} the identity on $\unsure{\mathfrak S_3|\mathcal G_3}$ \unsure{up to conjug\textsuperscript{acy}}))}

$\widetilde N_{0,3}$ denotes the group of \uline{outer} \unsure{automorphisms} of $\widehat{\mathcal G}_{0,3}^{+}$ which \deleted{\illegible{3}} \inserted{induce the identity in} $\unsure{\mathfrak S_3|\mathcal G_3}$. We note that
\[
\Gamma=\Gamma_{\mathbb Q}\hookrightarrow\widetilde\Gamma,
\]
and that the inverse image of $\Gamma_{\mathbb Q}$ in the extension (1) is none other than $M_{0,3}$ -- but for the moment I have no very plausible characterization to propose of $\Gamma_{\mathbb Q}$ as a subgroup of $\widetilde\Gamma$. We have a canonical splitting of the extension (1), via the subgroup $\widetilde\Gamma_{\sigma}$ of $M$ consisting of the \unsure{automorphisms} $u\in\widetilde M_{0,3}$ which commute with $\sigma_\infty$,



\scanpage{171}

so that $\widetilde M_{0,3}$ presents itself as a $\frac12$ direct product
\[
(3)\qquad \widetilde M_{0,3}\simeq\underbrace{\widehat{\mathcal G}_{0,3}^{+}}_{\mathrm{Sl}(2,\mathbb Z)'{}^\wedge}\cdot\widetilde\Gamma_\sigma.
\]

Let us consider the canonical homomorphism
\[
(4)\qquad \unsure{\widehat{\mathcal G}_{0,3}^{+}}\simeq\mathrm{Sl}(2,\mathbb Z)'{}^\wedge\xrightarrow{\varphi_0'}\mathrm{Sl}(2,\widehat{\mathbb Z})'\subset\mathrm{Gl}(2,\widehat{\mathbb Z})'
\]
and the \deleted{homomorphism}\inserted{map}
\[
\begin{aligned}
(5)\qquad \widetilde\Gamma_\sigma&\xrightarrow{\varphi_\sigma'}\mathrm{Gl}(2,\widehat{\mathbb Z})'\\
\underline u&\longmapsto\mu(t\sigma_\infty+u)\tau_\infty^\varepsilon=\mu\begin{pmatrix}u&\varepsilon t\\-t&\varepsilon u\end{pmatrix},
\end{aligned}
\]
where \deleted{$\unsure{t,u,\alpha=\alpha_2}$ \illegible{2}}\deleted{\unsure{$\varepsilon=\varepsilon_2$}}
\[
t,u\in\widehat{\mathbb Z},\quad\inserted{\mu\in\widehat{\mathbb Z}^{*},}\quad\varepsilon=\varepsilon_2\in\{\pm1\},\quad\alpha=\alpha_2\in\mathbb Z/2\mathbb Z
\]
are the quantities (depending on $\underline u\in\widetilde\Gamma_\sigma$) considered in the \S\ above, with
\[
\varepsilon=(-1)^\alpha.
\]
If need be, we shall write $t(\underline u)$, $\underline u(u)$, $\unsure{\alpha(\underline u)}$, $\varepsilon(\underline u)$ to indicate the dependence \unsure{with respect} to $\underline u$.

\deleted{\uline{Theorem}}\inserted{Let us consider the (continuous) map}
\[
(6)\qquad\varphi':\widetilde M_{0,3}\longrightarrow\mathrm{Gl}(2,\widehat{\mathbb Z})'
\]
defined by
\[
(7)\qquad\varphi'(a.\underline u)=\varphi_0'(a)\varphi_\sigma'(\underline u)\qquad\left\{\begin{aligned}&a\in\widehat{\mathcal G}_{0,3}^{+}\\&\underline u\in\widetilde\Gamma_\sigma\end{aligned}\right.
\]

\scanpage{172}
\mspage{385}

\uline{Theorem} The map $\varphi'$ is a group homomorphism.

This means, taking into account that $\varphi_0$ is already a group homomorphism, the theorem means that

a) $\varphi_\sigma'$ is a group homomorphism, i.e.
\[
(8)\qquad\varphi_\sigma'(\underline u'\underline u)=\varphi_\sigma'(\underline u')\unsure{\,\varphi_\sigma'(\underline u)}\qquad(\underline u,\underline u'\in\widetilde\Gamma_\sigma)
\]
b) we have, for $a\in\unsure{\mathcal G_{0,3}^{+}|\widehat{\mathcal G}_{0,3}^{+}}$, $\underline u\in\widetilde\Gamma_\sigma$
\[
(9)\qquad\varphi_0'\bigl(\underbrace{\underline u a\underline u^{-1}}_{\underline u(a)}\bigr)=\varphi_\sigma'(\underline u)\varphi_0'(a)\varphi_\sigma'\bigl(\underline u^{\deleted{\unsure{-1}}}\bigr)^{-1}.
\]

\deleted{To verify the}\inserted{But the} formula (9) is a fact \inserted{already established,} \deleted{\unsure{consequence of (8)}, \illegible{1} the} since the \deleted{\illegible{1}} very construction of the matrix $\varphi_\sigma'(\underline u)=\mu(t\sigma_\infty+u)\tau_\infty^{\unsure{\varepsilon}}$ in the preceding paragraphs was made in a way \deleted{to make} \unsure{such} that this formula should be \unsure{verified} (cf. (67)(68), and the corollary to theorem (77)\ldots). To verify (8), \deleted{we must} \inserted{we shall} first make explicit the quantities $g,h,f$ for a product of two operators $\unsure{u',u''|u',u}\in\widetilde M_{0,3}$, in terms of their values for $u',u$ themselves. An immediate calculation gives, in this sense, some ``cocycle relations''
\[
(10)\qquad\left\{\begin{aligned}
g(\underline u'\underline u)&=\unsure{\underline u'}(g(\underline u)).g(\underline u')\\
h(\underline u'\underline u)&=\underline u'(h(\underline u)).h(\underline u')\\
f(\underline u'\underline u)&=\unsure{f(\underline u')}.\underline u'(f(\underline u))
\end{aligned}\right.
\]

\scanpage{173}

\inserted{\uline{Remarks} 1)} We shall moreover have likewise, with the notations of the \S\ above
\[
(11)\qquad\left\{\begin{aligned}
\alpha(u'u)&=u'(\alpha(u)).\alpha(u')\\
\beta(u'u)&=u'(\beta(u)).\beta(u')
\end{aligned}\right.
\]

\deleted{\unsure{(N.B.}} 2) The formulas (10) are a little more \unsure{friendly|friendly} for $f$, than for $g$ and $h$, \inserted{$\alpha$ and $\beta$} (-- one would feel like changing notation yet again, and also replacing $g$ and $h$ by $g^{-1},h^{-1}$ (as previously \unsure{a $g_0$} had become $f^{-1}$\ldots).

Let us again denote by the same letters $g(u),h(u)$ etc. the images of these quantities in $\mathrm{Sl}(2,\widehat{\mathbb Z})'$, \deleted{the formulas \unsure{(10)} \inserted{(11)} \unsure{give}} the analogous formulas in $\mathrm{Sl}(2,\widehat{\mathbb Z})'$ (or $(\mathrm{Gl}(2,\widehat{\mathbb Z})^{\pm1})'$ for \deleted{\unsure{(11)}} the $\alpha$ of (11)) which are written
\[
(12)\qquad\left\{\begin{aligned}
g(\underline u'\underline u)&=\mathrm{int}(\varphi_\sigma'(\underline u'))(g(\underline u)).g(\underline u')\\
h(\underline u'\underline u)&=\mathrm{int}(\varphi_\sigma'(\underline u'))(h(\underline u)).h(\underline u')\\
f(\underline u'\underline u)&=\deleted{\unsure{\mathrm{int}(\varphi_\sigma'(\underline u'))}}f(\underline u').\bigl(\mathrm{int}\unsure{(\varphi_\sigma(\underline u'))}(f(\underline u))\bigr)\\
\inserted{\alpha(\underline u'\underline u)}&\inserted{{}=\mathrm{int}\,\varphi_\sigma'(\underline u')(\alpha(\underline u)).\alpha(\underline u')}\\
\inserted{\beta(\underline u'\underline u)}&\inserted{{}=\mathrm{int}\,\varphi_\sigma'(\underline u')(\beta(\underline u)).\beta(\underline u')}
\end{aligned}\right.
\]
in $\mathrm{Gl}(2,\widehat{\mathbb Z})'$.

\deleted{To verify}\inserted{These formulas} in $\mathrm{Gl}(2,\widehat{\mathbb Z})'$ do indeed have an entirely explicit character from the point of view of effective computation, unlike

\scanpage{174}
\mspage{386}

the case of the initial formulas in $\widehat{\Pi}_{0,3}$ or $\unsure{\widehat{\widetilde{\Pi}}_{0,3}}$, where quantities like $\underline{u}'(\underbrace{g(\underline{u})}_{g})$ cannot be made explicit \deleted{\illegible{1}} \unsure{by} simple algebraic formulas in terms of the data \inserted{$\unsure{g,h,f,}$} $g',h',f'$ \deleted{\unsure{resp.}} $=g(\underline{u})$ etc. which determine $\underline{u}$ and $\underline{u}'$.

These formulas are supplemented, for the record, by
\[
\tag{13}
\mu(\underline{u}'\underline{u})=\mu(\underline{u}')\mu(\underline{u})
\]
implying in particular
\[
\tag{14}
\varepsilon_2(\underline{u}'\underline{u})=\varepsilon_2(\underline{u}')\varepsilon_2(\underline{u}).
\]

To carry out the calculation of $\varphi_{\sigma}(\underline{u}'\mathbin{.}\underline{u})$ and verify that it is equal to $\varphi_{\sigma}(\underline{u}')\varphi_{\sigma}(\underline{u})$, we \deleted{\unsure{shall have}} shall introduce the coefficients $r,s,t,u$ of $f=f(\underline{u})$, $r',s',t',u'$ of \inserted{$f'=$} $f(\underline{u}')$, and $r'',s'',t'',u''$ of $f''=f(\underline{u}'\underline{u})$, likewise \inserted{$\mu,\mu',\mu''$,} $\varepsilon,\varepsilon',\varepsilon''$, whence
\[
\tag{15}
\deleted{\text{and }\unsure{\text{relations}}}\qquad
\mu''=\mu\mu',\quad \varepsilon''=\varepsilon\varepsilon'
\]
and the formula in $f$ in (12) is therefore written
\[
\tag{16}
f''=f'\underline{u}'f\underline{u}'^{-1}
\]
where we have written, as in the \S\ above, $\underline{u}'$ instead of $\varphi_{\sigma}(\underline{u}')$, \deleted{\unsure{In}} In explicit terms
\[
\begin{pmatrix}r''&s''\\t''&u''\end{pmatrix}
=\begin{pmatrix}r'&s'\\t'&u'\end{pmatrix}
\unsure{\mu'}\begin{pmatrix}u'&\varepsilon't'\\-t'&\varepsilon'u'\end{pmatrix}
\begin{pmatrix}r&s\\t&u\end{pmatrix}
\unsure{\deleted{\frac{\mu'}{\mu'}}|\frac{\mu'}{\mu'}}
\begin{pmatrix}\varepsilon'u'&-\varepsilon't'\\t'&u'\end{pmatrix}
\]
\margin{N.B. $\det\underline{u}'=\mu'$}